% Generated draft: requires mathematical and source review.
\documentclass[11pt]{article}
\usepackage[margin=25mm]{geometry}
\usepackage{fontspec}
\setmainfont{Latin Modern Roman}
\usepackage{amsmath,amssymb,mathtools,unicode-math}
\setmathfont{Latin Modern Math}
\usepackage{enumitem,needspace,xurl,hyperref}
\setlength{\parindent}{0pt}
\setlength{\parskip}{5pt}
\setlength{\emergencystretch}{3em}
\newcommand{\N}{\mathbb{N}}
\newcommand{\Z}{\mathbb{Z}}
\newcommand{\Q}{\mathbb{Q}}
\newcommand{\R}{\mathbb{R}}
\newcommand{\C}{\mathbb{C}}
\newcommand{\F}{\mathbb{F}}
\newcommand{\A}{\mathbb{A}}
\newcommand{\PP}{\mathbb{P}}
\newcommand{\E}{\mathbb{E}}
\newcommand{\eps}{\varepsilon}
\newcommand{\dd}{\,\mathrm{d}}
\begin{document}
\section*{UnsolvedMath ID: 801125}
\textbf{Source problem code:} catalogue entry
\subsection{Definitions and mathematical statement}
A type \(\mathrm{II}_1\) factor is a von Neumann algebra \(N\) with center \(\C 1\), no nonzero minimal projections, and a faithful normal tracial state \(\tau\) satisfying \(\tau(1)=1\). No separability assumption is stated in the supplied record. Write \(\mathfrak c=2^{\aleph_0}\) for the cardinality of the continuum.

The supplied record does not define the invariant \(h(N)\) or the term ``sequential commutation orbit.'' Consequently, the following is a conditional formulation requiring those definitions from the source. Let \(\mathscr O_{\mathrm{sc}}(N)\) denote the family of sequential commutation orbits intended by the source, assuming these are subsets of \(N\). For a subset \(S\subseteq N\), let \(W_N^*(S)\) denote the smallest unital von Neumann subalgebra of \(N\) containing \(S\).

Interpreting ``continuum many'' as exactly \(\mathfrak c\), the two questions are:
\begin{enumerate}
\item Does there exist a type \(\mathrm{II}_1\) factor \(N\) such that
\[
h(N)=0
\quad\text{and}\quad
\lvert\mathscr O_{\mathrm{sc}}(N)\rvert=\mathfrak c?
\]
\item Does there exist a type \(\mathrm{II}_1\) factor \(N\) such that
\[
h(N)=0
\quad\text{and}\quad
\forall O\in\mathscr O_{\mathrm{sc}}(N),\qquad W_N^*(O)\ne N?
\]
\end{enumerate}
The existential witnesses in the two questions need not coincide. The meaning of \(h\), the underlying elements and ambient algebra used to define the orbits, and any source-specific hypotheses remain unresolved; this formulation does not supply those missing definitions.
\subsection{Short English statement}
The source asks two separate existence questions: can a type \(\mathrm{II}_1\) factor with \(h(N)=0\) have continuum many sequential commutation orbits, and can such a factor fail to be generated by any single such orbit? The invariant and orbit terminology require source verification before these questions become fully specified mathematical statements.
\subsection{Sources}
\begin{itemize}
\item Local UnsolvedMath registry: permanent ID \(801125\), exact source problem code \texttt{catalogue entry}, title ``Ultraproduct \(\mathrm{II}_1\) Factors --- Problem 26.'' The supplied record attributes the problem to Elayavalli, with proposed year 2025, and labels it open; this status is supplied metadata, not an independently verified mathematical assessment.
\item Supplied source list: Elayavalli, \emph{50 Open Problems: Ultraproduct \(\mathrm{II}_1\) Factors} (2025), Problem 26. Supplied URL: \url{https://arxiv.org/abs/2511.20377}. The record reports access on 2026-10-16 and extraction by \texttt{arxiv\_latex}; these are supplied provenance details, not a claim of independent access or verification.
\end{itemize}
\end{document}
